% ============================================================
%  preamble.tex — Moteur de style du cours
%
%  Ne modifiez pas ce fichier pour personnaliser votre cours :
%  utilisez config.tex.
%
%  Compilation : lualatex main.tex (2 passes)
% ============================================================

% ------------------------------------------------------------
% Mise en page et polices
% ------------------------------------------------------------
\usepackage[a4paper,margin=2cm,headheight=19pt]{geometry}
\usepackage{fontspec}
\setmainfont{Liberation Sans}
\setsansfont{Liberation Sans}
\setmonofont{JetBrains Mono}[
  Path=fonts/,
  UprightFont=JetBrainsMono-Regular.ttf,
  BoldFont=JetBrainsMono-Bold.ttf,
  ItalicFont=JetBrainsMono-Italic.ttf,
  BoldItalicFont=JetBrainsMono-BoldItalic.ttf,
  Scale=0.95
]

\usepackage[french]{babel}
\usepackage{microtype}
\usepackage{amsmath,amssymb}
\usepackage{graphicx}
\usepackage{float}
\usepackage{xcolor}
\usepackage{tikz}
\usetikzlibrary{arrows.meta,positioning,shapes.geometric,calc,fit,backgrounds}
\usepackage[most]{tcolorbox}
\tcbuselibrary{listings,skins,breakable}
\usepackage{listings}
\usepackage{enumitem}
\usepackage{booktabs}
\usepackage{tabularx}
\usepackage{array}
\usepackage{fancyhdr}
\usepackage{lastpage}
\usepackage{caption}
\usepackage{hyperref}

% ------------------------------------------------------------
% Couleurs
% ------------------------------------------------------------
\definecolor{deepblue}{HTML}{0F2A4A}
\definecolor{teal}{HTML}{0F766E}
\definecolor{warm}{HTML}{C2410C}
\definecolor{codebg}{HTML}{F6F8FA}
\definecolor{codegray}{HTML}{6B7280}
\definecolor{softblue}{HTML}{EAF2F8}
\definecolor{softgray}{HTML}{F4F6F8}
\definecolor{softteal}{HTML}{E6F4F1}

\hypersetup{
  colorlinks=true,
  linkcolor=deepblue,
  urlcolor=teal,
  citecolor=teal,
  pdftitle={\coursTitre\ — \coursSousTitre},
  pdfauthor={\coursAuteur},
  pdfsubject={\coursSousTitre},
  pdfkeywords={cours, support de cours, formation}
}

% ------------------------------------------------------------
% Typographie générale
% ------------------------------------------------------------
\setlength{\parindent}{0pt}
\setlength{\parskip}{0.45em}
\setlength{\emergencystretch}{2em}
\raggedbottom
\renewcommand{\arraystretch}{1.2}
\setlist[itemize]{leftmargin=1.4em,itemsep=2pt,topsep=3pt}
\setlist[enumerate]{leftmargin=1.7em,itemsep=2pt,topsep=3pt}

\captionsetup{font={small,sf},labelfont={bf,color=deepblue}}
\renewcommand{\lstlistingname}{Extrait de code}
\renewcommand{\lstlistlistingname}{Liste des extraits de code}

% ------------------------------------------------------------
% En-têtes et pieds de page
% ------------------------------------------------------------
\pagestyle{fancy}
\fancyhf{}
\fancyhead[L]{\sffamily\footnotesize\color{codegray}\coursEntete}
\fancyhead[R]{\sffamily\footnotesize\color{codegray}\leftmark}
\fancyfoot[C]{\sffamily\footnotesize\color{codegray}\thepage\ / \pageref{LastPage}}
\renewcommand{\headrulewidth}{0.4pt}
\renewcommand{\headrule}{\hbox to\headwidth{\color{teal!30}\leaders\hrule height \headrulewidth\hfill}}

% ------------------------------------------------------------
% Chapitres et sections personnalisés
% ------------------------------------------------------------
\newcounter{chapter}
\renewcommand{\thechapter}{\arabic{chapter}}

% \chapitre{numéro}{Titre}
% Le numéro est affiché. Il doit suivre l'ordre des chapitres.
\newcommand{\chapitre}[2]{%
  \clearpage
  \setcounter{chapter}{#1}%
  \phantomsection
  \pdfbookmark[1]{Chapitre #1\ — #2}{chapter#1}%
  \addcontentsline{toc}{section}{\protect\textbf{Chapitre #1\ — #2}}%
  \markboth{Chapitre #1\ — #2}{}%
  \thispagestyle{plain}%
  \begin{tikzpicture}[remember picture,overlay]
    \fill[deepblue] (current page.north west) rectangle
      ([yshift=-2.7cm]current page.north east);
    \fill[teal] ([yshift=-2.62cm]current page.north west) rectangle
      ([yshift=-2.7cm]current page.north east);
    \node[anchor=west,text=white,font=\sffamily\bfseries\fontsize{30}{34}\selectfont]
      at ([xshift=2cm,yshift=-1.15cm]current page.north west) {Chapitre #1};
    \node[anchor=west,text=teal!50,font=\sffamily\bfseries\fontsize{10}{12}\selectfont]
      at ([xshift=2cm,yshift=-2.05cm]current page.north west)
      {\addfontfeatures{LetterSpace=12}\coursBandeauChapitre};
  \end{tikzpicture}%
  \par\vspace*{0.7cm}
  {\sffamily\bfseries\color{deepblue}\fontsize{22}{26}\selectfont #2\par}
  \vspace{0.25cm}
  {\color{teal}\rule{\linewidth}{1.2pt}}
  \vspace{0.65cm}
}

% \soussection{numéro}{Titre}
\newcommand{\soussection}[2]{%
  \phantomsection
  \pdfbookmark[2]{#1 — #2}{sec:#1}%
  \addcontentsline{toc}{subsection}{\protect\numberline{#1}#2}%
  \markright{#1 — #2}%
  {\sffamily\bfseries\color{teal}\large #1\quad #2\par}
  \vspace{0.25cm}
}

% \soussoussection{numéro}{Titre}
\newcommand{\soussoussection}[2]{%
  \phantomsection
  \pdfbookmark[3]{#1 — #2}{subsec:#1}%
  \addcontentsline{toc}{subsubsection}{\protect\numberline{#1}#2}%
  {\sffamily\bfseries\color{deepblue}\normalsize #1\quad #2\par}
  \vspace{0.15cm}
}

% ------------------------------------------------------------
% Encadrés graphiques
% Tous s'utilisent de la même façon :
%   \begin{nom}{} ... \end{nom}
% ------------------------------------------------------------
\newtcolorbox{objectifs}[1][]{
  enhanced,breakable,
  title={\sffamily\bfseries Objectifs du chapitre},
  colback=deepblue!2,colframe=deepblue!35,
  colbacktitle=deepblue!8,coltitle=deepblue,
  fonttitle=\small,boxrule=0.5pt,arc=1.5mm,
  left=8pt,right=8pt,top=7pt,bottom=7pt,#1
}

\newtcolorbox{definition}[1][]{
  enhanced,breakable,
  title={\sffamily\bfseries Définition},
  colback=white,colframe=deepblue!30,
  colbacktitle=deepblue!6,coltitle=deepblue,
  fonttitle=\small,boxrule=0.5pt,arc=1.5mm,
  left=8pt,right=8pt,top=7pt,bottom=7pt,#1
}

\newtcolorbox{retenir}[1][]{
  enhanced,breakable,
  title={\sffamily\bfseries À retenir},
  colback=white,colframe=teal!45,
  colbacktitle=teal!8,coltitle=teal!80!black,
  fonttitle=\small,boxrule=0.5pt,arc=1.5mm,
  left=8pt,right=8pt,top=7pt,bottom=7pt,#1
}

\newtcolorbox{attention}[1][]{
  enhanced,breakable,
  title={\sffamily\bfseries Attention},
  colback=red!3,colframe=red!45!black,
  colbacktitle=red!8,coltitle=red!60!black,
  fonttitle=\small,boxrule=0.5pt,arc=1.5mm,
  left=8pt,right=8pt,top=7pt,bottom=7pt,#1
}

\newtcolorbox{piege}[1][]{
  enhanced,breakable,
  title={\sffamily\bfseries Piège classique},
  colback=warm!3,colframe=warm!45,
  colbacktitle=warm!9,coltitle=warm!80!black,
  fonttitle=\small,boxrule=0.5pt,arc=1.5mm,
  left=8pt,right=8pt,top=7pt,bottom=7pt,#1
}

\newtcolorbox{conseil}[1][]{
  enhanced,breakable,
  title={\sffamily\bfseries Conseil},
  colback=teal!3,colframe=teal!45,
  colbacktitle=teal!10,coltitle=teal!80!black,
  fonttitle=\small,boxrule=0.5pt,arc=1.5mm,
  left=8pt,right=8pt,top=7pt,bottom=7pt,#1
}

\newtcolorbox{choix}[1][]{
  enhanced,breakable,
  title={\sffamily\bfseries Choix recommandé},
  colback=softblue,colframe=deepblue!40,
  colbacktitle=deepblue!8,coltitle=deepblue,
  fonttitle=\small,boxrule=0.5pt,arc=1.5mm,
  left=8pt,right=8pt,top=7pt,bottom=7pt,#1
}

\newtcolorbox{atelier}[1][]{
  enhanced,breakable,
  title={\sffamily\bfseries Atelier},
  colback=softgray,colframe=codegray!45,
  colbacktitle=codegray!12,coltitle=black!75,
  fonttitle=\small,boxrule=0.5pt,arc=1.5mm,
  left=8pt,right=8pt,top=7pt,bottom=7pt,#1
}

\newtcolorbox{resume}[1][]{
  enhanced,breakable,
  title={\sffamily\bfseries En résumé},
  colback=softteal,colframe=teal!25,
  colbacktitle=teal!8,coltitle=teal!80!black,
  fonttitle=\small,boxrule=0.4pt,arc=1.2mm,
  left=8pt,right=8pt,top=6pt,bottom=6pt,#1
}

% ------------------------------------------------------------
% Langages et styles de code
% ------------------------------------------------------------
\lstdefinelanguage{Java25}{
  language=Java,
  morekeywords={
    record, sealed, permits, non-sealed, yield, var, when
  },
  morekeywords=[2]{
    @SpringBootApplication, @RestController, @Controller, @Service, @Repository,
    @Component, @Configuration, @Bean, @Autowired, @Value, @ConfigurationProperties,
    @GetMapping, @PostMapping, @PutMapping, @DeleteMapping, @PatchMapping,
    @RequestBody, @PathVariable, @RequestParam, @ResponseStatus,
    @RestControllerAdvice, @ControllerAdvice, @ExceptionHandler,
    @CrossOrigin, @PostConstruct, @PreDestroy, @Valid, @NotNull, @NotBlank,
    @Size, @Min, @Max, @Tag, @Operation, @ApiResponse, @Parameter, @Schema
  }
}

\lstdefinelanguage{GradleGroovy}{
  keywords={
    plugins, id, version, repositories, mavenCentral, dependencies,
    implementation, testImplementation, runtimeOnly, compileOnly,
    java, toolchain, languageVersion, JavaLanguageVersion,
    group, sourceCompatibility, tasks, withType, named, test
  }
}

\lstdefinelanguage{PropertiesLang}{
  keywords={},
  sensitive=false
}

% Style commun à tous les blocs de code
\lstdefinestyle{coursbase}{
  backgroundcolor=\color{codebg},
  basicstyle=\ttfamily\bfseries\coursCodeSize,
  keywordstyle=\color{deepblue}\bfseries,
  keywordstyle=[2]\color{teal}\bfseries,
  commentstyle=\color{teal!70!black}\itshape,
  stringstyle=\color{warm!80!black},
  numbers=none,
  frame=single,rulecolor=\color{deepblue!20},
  breaklines=true,breakatwhitespace=false,
  postbreak=\mbox{\textcolor{teal}{$\hookrightarrow$}\space},
  showstringspaces=false,tabsize=2,columns=fullflexible,
  keepspaces=true,upquote=true,
  captionpos=b,xleftmargin=0.5em,framexleftmargin=0.3em,
  lineskip=-0.6pt,
  aboveskip=4pt,belowskip=4pt
}

\lstdefinestyle{java}{style=coursbase,language=Java25}
\lstdefinestyle{gradle}{style=coursbase,language=GradleGroovy}
\lstdefinestyle{properties}{style=coursbase,language=PropertiesLang}
\lstdefinestyle{bash}{style=coursbase,language=bash}
\lstdefinestyle{json}{style=coursbase}
\lstdefinestyle{sql}{style=coursbase,language=SQL}
\lstdefinestyle{python}{style=coursbase,language=Python}
\lstdefinestyle{html}{style=coursbase,language=HTML}
\lstdefinestyle{xml}{style=coursbase,language=XML}
\lstdefinestyle{c}{style=coursbase,language=C}
\lstdefinestyle{cpp}{style=coursbase,language=C++}
\lstdefinestyle{text}{style=coursbase}

% ------------------------------------------------------------
% Raccourcis de mise en forme
% ------------------------------------------------------------
% Code court dans le texte : \code{nom_de_variable}
\DeclareRobustCommand{\code}[1]{%
  \texorpdfstring{\colorbox{codebg}{\texttt{\coursCodeSize\bfseries\color{deepblue!90}\detokenize{#1}}}}{\detokenize{#1}}%
}

% Terme important : \motcle{concept}
\newcommand{\motcle}[1]{\textsf{\color{teal!80!black}#1}}

% Alias de compatibilité avec l'ancien template
\newcommand{\spring}[1]{\motcle{#1}}

% ------------------------------------------------------------
% Table des matières : largeur des numéros
% ------------------------------------------------------------
\makeatletter
\renewcommand*\l@section{\@dottedtocline{1}{0.5em}{2.8em}}
\renewcommand*\l@subsection{\@dottedtocline{2}{2.0em}{3.4em}}
\renewcommand*\l@subsubsection{\@dottedtocline{3}{4.2em}{4.2em}}
\makeatother
